%% Đề 8 -- de-giua-hoc-ki-2-toan-9-nam-2022-2023-phong-gddt-bac-son-lang-son.pdf
\deheader{ĐỀ SỐ 8}%
  {Đề kiểm tra Giữa Học kỳ — Năm học 2022--2023}%
  {Phòng GD\&ĐT | Chỉ phần Đại số}
\addcontentsline{toc}{section}{Đề 8 — Đề kiểm tra Giữa HK — Năm học 2022--2023}

\begin{de_cau}{Câu 1}{2,0 điểm}
\begin{enumerate}[label=\alph*)]
\item Tính: $\dfrac{1}{2}\sqrt{48}-\dfrac{3}{2}\sqrt{12}+\sqrt{75}$.
\item Cho $P=\dfrac{x+2}{x\sqrt{x}-1}+\dfrac{\sqrt{x}+1}{x+\sqrt{x}+1}-\dfrac{1}{\sqrt{x}-1}$, $x>0, x\neq 1$.
Rút gọn $P$.
\end{enumerate}
\end{de_cau}

\begin{loi_giai_cau}
\textbf{a)} $\frac{1}{2}\cdot 4\sqrt{3} - \frac{3}{2}\cdot 2\sqrt{3} + 5\sqrt{3} = 2\sqrt{3}-3\sqrt{3}+5\sqrt{3} = 4\sqrt{3}$.

\textbf{b)} Mẫu chung $(\sqrt{x}-1)(x+\sqrt{x}+1)$:
\[P = \frac{x+2+(\sqrt{x}+1)(\sqrt{x}-1)-(x+\sqrt{x}+1)}{(\sqrt{x}-1)(x+\sqrt{x}+1)}
= \frac{x+2+x-1-x-\sqrt{x}-1}{\ldots} = \frac{x-\sqrt{x}}{(\sqrt{x}-1)(x+\sqrt{x}+1)}.\]
\[= \frac{\sqrt{x}(\sqrt{x}-1)}{(\sqrt{x}-1)(x+\sqrt{x}+1)} = \frac{\sqrt{x}}{x+\sqrt{x}+1}.\]
\end{loi_giai_cau}

\begin{de_cau}{Câu 2}{2,0 điểm}
\begin{enumerate}[label=\alph*)]
\item Vẽ đồ thị hàm số $y = \dfrac{1}{2}x^2$.
\item Tìm tọa độ giao điểm của $(P)\colon y=x^2$ và $(d)\colon y=x+2$.
\end{enumerate}
\end{de_cau}

\begin{loi_giai_cau}
\textbf{a)} Bảng giá trị $y=\frac{1}{2}x^2$:
\begin{center}
\begin{tabular}{c|ccccc}$x$&$-2$&$-1$&$0$&$1$&$2$\\\hline$y$&$2$&$\frac{1}{2}$&$0$&$\frac{1}{2}$&$2$\end{tabular}
\end{center}

\textbf{b)} Phương trình hoành độ giao điểm: $x^2=x+2\Rightarrow x^2-x-2=0\Rightarrow(x-2)(x+1)=0$.

$x=2\Rightarrow y=4$: điểm $A(2;4)$.
$x=-1\Rightarrow y=1$: điểm $B(-1;1)$.
\end{loi_giai_cau}

\begin{de_cau}{Câu 3}{1,5 điểm}
Một ô tô đi từ A đến B với vận tốc 60 km/h. Lúc về ô tô đi nhanh hơn lúc đi 20 km/h nên về sớm hơn
đi 1 giờ. Tính quãng đường AB.
\end{de_cau}

\begin{loi_giai_cau}
Gọi $AB = d$ (km). Vận tốc đi: 60 km/h, về: 80 km/h.

\[\frac{d}{60}-\frac{d}{80}=1\Rightarrow\frac{4d-3d}{240}=1\Rightarrow d=240\text{ km}.\]

Vậy $AB = \boxed{240}$ km.
\end{loi_giai_cau}

\begin{de_cau}{Câu 4}{1,0 điểm}
Cho phương trình $x^2-2(m-2)x+m^2-3m=0$ (tham số $m$). Tìm $m$ để phương trình có hai nghiệm
$x_1, x_2$ thỏa mãn $x_1^2+x_2^2\leq 10$.
\end{de_cau}

\begin{loi_giai_cau}
$\Delta'=(m-2)^2-(m^2-3m) = m^2-4m+4-m^2+3m = -m+4$.

Để có hai nghiệm: $\Delta'\geq 0\Rightarrow m\leq 4$.

Vi-ét: $S=x_1+x_2=2(m-2)$, $P=x_1x_2=m^2-3m$.

$x_1^2+x_2^2=S^2-2P=4(m-2)^2-2(m^2-3m)=2m^2-10m+16$.

$2m^2-10m+16\leq 10\Rightarrow 2m^2-10m+6\leq 0\Rightarrow m^2-5m+3\leq 0$.

$m=\dfrac{5\pm\sqrt{13}}{2}$, nên $\dfrac{5-\sqrt{13}}{2}\leq m\leq \dfrac{5+\sqrt{13}}{2}$.

Kết hợp $m\leq 4$: $\boxed{\dfrac{5-\sqrt{13}}{2}\leq m\leq 4}$.
\end{loi_giai_cau}
